\section{Codabench submissions, bundles and release}\label{app:K}

\paragraph{Submissions.} Every submission under our Codabench account (user \texttt{dan23rr}) in both phases, read from the Codabench API on 28 October 2026. Accuracy is given as items correct over items; a parent submission sends one bundle to both tasks, with one child per task it was evaluated on (the parents of 13 September also carry the cancelled children listed below the table).

{\footnotesize\setlength{\tabcolsep}{4pt}

\begin{longtable}{L{1.3cm}L{1.9cm}L{1.7cm}L{1.5cm}L{5.9cm}}

\toprule

\textbf{ID} & \textbf{Task} & \textbf{Date (UTC)} & \textbf{Accuracy} & \textbf{What it is} \\

\midrule

\endfirsthead

\toprule

\textbf{ID} & \textbf{Task} & \textbf{Date (UTC)} & \textbf{Accuracy} & \textbf{What it is} \\

\midrule

\endhead

\bottomrule

\endlastfoot

907310 & warmup & 2026-08-29 & 8/17 & Null probe, every item non-robust; floor pre-registered at 8/17 (full report, Section 3.1). \\

907677 & warmup & 2026-08-29 & 9/17 & Official baseline, unchanged, instrumented to log its runtime (Appendix D). \\

907967 & warmup & 2026-08-29 & 9/17 & Early version of the measuring entry, with a time budget, as its released bundle documents; its logs give the container's GPU and library versions (Appendix D). \\

909567 & warmup & 2026-08-31 & 9/17 & Later version of the same entry, as its released bundle documents. \\

909843 & warmup & 2026-08-31 & 9/17 & V4, the allocation entry, submitted once as a measurement against its pre-registered rule (Appendix G). \\

909853 & warmup & 2026-08-31 & 10/17 & v5: every model measured by itself (H16, Appendix A). \\

909976 & warmup & 2026-08-31 & 11/17 & v6: H16 with a gpt-oss threshold calibrated on \texttt{val-sample} (H17); with the next three, the four executions of H18. \\

909994 & warmup & 2026-08-31 & 9/17 & v6, re-executed (H17, H18). \\

910002 & warmup & 2026-08-31 & 11/17 & v6, re-executed (H17, H18). \\

910013 & warmup & 2026-08-31 & 11/17 & v6, re-executed (H17, H18). \\

912096 & Main (original) & 2026-09-02 & 11/21 & v6, the measuring entry, on the original Main task (Section 6). \\

912097 & Small (original) & 2026-09-02 & 10/19 & v6 on the original Small task (Section 6). \\

924968 & both tasks (parent) & 2026-09-13 & parent & V8 Main bundle, sent to both original tasks: parent of 924969 and 924970. \\

924969 & Main (original) & 2026-09-13 & 11/21 & V8 on the original Main task (Section 6). \\

924970 & Small (original) & 2026-09-13 & coverage 0 & V8 Main bundle on the Small task, rejected before import because it lacks \texttt{small.txt} (coverage 0). \\

924972 & both tasks (parent) & 2026-09-13 & parent & V8 Small bundle, sent to both original tasks: parent of 924973 and 924974. \\

924973 & Main (original) & 2026-09-13 & coverage 0 & V8 Small bundle on the Main task, rejected before import because of \texttt{small.txt} (coverage 0). \\

924974 & Small (original) & 2026-09-13 & 10/19 & V8 on the original Small task (Section 6). \\

925822 & Main (original) & 2026-09-14 & 11/21 & V8 Main bundle again, Main task only, uploaded by the AI assistant without a specific instruction (AI usage statement). \\

936802 & Main & 2026-09-21 & 10/21 & V8 on the replacement Main task (Section 6, Appendix B). \\

936806 & Small & 2026-09-21 & 7/14 & V8 on the replacement Small task (Section 6, Appendix B). \\

1016713 & Main & 2026-10-28 & 10/21 & V8, unchanged, on the current validation task, after the change of 15 October (Section 6); nominated for the final test, Main track. \\

1016730 & Small & 2026-10-28 & 7/14 & V8, unchanged, on the current validation task, after the change of 15 October (Section 6). \\

1016732 & Main & 2026-10-28 & 10/21 & All-robust constant, control (Section 6); nominated for the final test, Main track. \\

1016733 & both tasks (parent) & 2026-10-28 & parent & Constant Small bundle: parent of 1016734 and 1016735. \\

1016734 & Main & 2026-10-28 & coverage 0 & Constant Small bundle sent to the Main task in error (AI usage statement); rejected before import (``small.txt is not allowed for a main-track submission''), as 924973 was on the original Main task. \\

1016735 & Small & 2026-10-28 & 7/14 & All-robust constant, control (Section 6). \\

\end{longtable}

}

Also under our account, and not described in this report: 928172, 928173, 928174, 928175, 929013 and 929014, children of 924968 and 924972 created on 16 September and cancelled, without scores.

\paragraph{Bundles.} Our entry on the current tasks is V8, with the all-robust constant as control; we nominated 1016713 and 1016732 for the final test, on the Main track. Neither loads a model or reads data at run time. V8's two parameters were fixed on the six public datasets, which are its training data; Appendix D (full report, Section 13) gives them, the scripts that recompute them and the instructions to reproduce its score. SHA-256 of each bundle:

\begin{itemize}

\item \texttt{solution-v8-main.zip}: V8: a single \texttt{solution.py} defining \texttt{are\_robust(model\_id, problems)}, which returns \texttt{[len(p) >= 98 for p in problems]}. Select only the Main task.\\
{\footnotesize\texttt{7b70f15c51752dcf4285e477d58c541d0ecf052d833c325c0702dbbd054ffb5f}}

\item \texttt{solution-v8-small.zip}: The same \texttt{solution.py}, with the empty marker file \texttt{small.txt} that the Small task requires. Select only the Small task.\\
{\footnotesize\texttt{0b83b55f94dc16b9852a7c2b35d71ac5787e43d81afd6396dee00fb5e011e500}}

\item \texttt{solution-costante-main.zip}: The all-robust constant: a single \texttt{solution.py} whose \texttt{are\_robust} returns \texttt{True} for every problem, written to accept the argument orders of both official descriptions of the interface and the legacy two-argument call. Select only the Main task.\\
{\footnotesize\texttt{35c47bbef717a5d4c5b6c7786bddbab9d87a02483e4f417cd34fd4ecf669c67d}}

\item \texttt{solution-costante-small.zip}: The same \texttt{solution.py}, with \texttt{small.txt}. Select only the Small task.\\
{\footnotesize\texttt{f798c2979b136a84c22e6d994bf0228bbfeffa044ac68f7868e13859b4c6075a}}

\end{itemize}

The V8 docstring, frozen with the bundle on 13 September, contains three statements that this report corrects: it calls 98 the median of the character counts (it is the lower median) and gives the rule's balanced accuracy as 0.658, the figure for a threshold of 99 (Appendix C, item 37); and it says that the official score changes sign between identical runs, quoting AUROC 0.513 and 0.398 with robust as the positive class on the raw score, where Section 4 gives the same runs as 0.487 and 0.602 in the regressor's own orientation and finds the difference not significant.

\paragraph{Reproducing our score.} Upload each bundle with only its own task selected: the Main bundle on task 36237, the Small bundle on task 36241. No model and no data are needed. On the current validation tasks the expected results are those of 1016713 and 1016730, 10/21 and 7/14; Appendix D gives the instructions and expected results of 23 September, before the change of 15 October.

\paragraph{Release.} Every pre-registration, outcome file and analysis script, the extraction outputs, the submission bundles as submitted, our perturbed variants of the public problems, our snapshots of the Codabench leaderboard and the full report (draft v2.4) are released with this report at https://github.com/Dan23RR/aimo-instrument, which keeps the layout of our working folder so that the file paths cited here are valid there, with other participants' names, usernames and passages about their entries redacted. The organizers' getting-started files inside our bundles and our variants are released with their permission (Appendix H); copies of the organizers' datasets are not, and in our variant files the rows of the original problems keep their identifiers but not their texts, while five original problems remain, as submitted, in the calibration file of the v6 bundles. Our correspondence, internal notes and the drafts before v2.4 that Appendix D lists are not released. On 7 September we asked the organizers whether measurements taken on their virtual machine may appear in this report and undertook to share our results with them, and we repeated the question on 2 November; the same day they answered that this is acceptable, and the results reach them with this report and the release. Deviations from the pre-registrations are listed in Appendix D.
